%%% Graphoid Axioms
\section{Graphoid Axioms}\label{apx:graphoid}
\begin{definition}[Graphoid Axioms]\label{def:graphoid}
  For disjoint subsets $\rmX,\rmY,\rmZ,\rmW\subseteq\verts$, a
  semi-graphoid independence model $\indeps{\cdot}$ obeys the following
  properties~\citep{dawid1979,pearl1987}:
  \begin{enumerate}
  \item Symmetry:
    $\rmX \indep{} \rmY \mid \rmZ \implies \rmY\indep{}\rmX\mid\rmZ$
  \item Decomposition:
    $\rmX \indep{} \rmY \cup \rmW \mid \rmZ \implies
    \rmX\indep{}\rmY\mid\rmZ$
  \item Weak Union:
    $\rmX\indep{}\rmY\cup\rmW\mid\rmZ \implies \rmX\indep{}\rmY \mid
    \rmW\cup\rmZ$
  \item Contraction:
    $(\rmX\indep{}\rmY \mid \rmW\cup\rmZ) \wedge
    (\rmX\indep{}\rmW\mid\rmZ)\implies \rmX\indep{}\rmY\cup\rmW\mid\rmZ$
  \end{enumerate}
  The independence model is called a graphoid if it also obeys the
  following additional property~\citep{pearl1985}:
  \begin{enumerate}[resume]
  \item Intersection:
    $(\rmX \indep{} \rmY \mid \rmZ \cup \rmW) \wedge
    (\rmX \indep{} \rmW \mid \rmZ \cup \rmY)
    \implies \rmX \indep{} \rmY \cup \rmW \mid
    \rmZ$
  \end{enumerate}
  All graph independence models $\indeps{\gr}$ are guaranteed to
  obey all the graphoid axioms~\citep{lauritzen2018}. However,
  probabilistic independence models $\indeps{\pr}$ are only guaranteed to
  obey the semi-graphoid axioms. However if the joint distribution is
  strictly positive, then $\indeps{\pr}$ is guaranteed to be a
  graphoid as well~\citep{pearl1988,pearl1989}.
\end{definition}

%%% Proofs
\section{Proofs}\label{apx:proofs}
%%%% k dep to assoc
\kdeptoassoc*
\begin{proof}
  Consider the family of association witnesses contained in $\rmZ$,
  \begin{equation*}
    \mathcal{W} \coloneq
    \{\rmW \subseteq \rmZ : \kassoc{\rmS}{Y,X}{\rmW}\}.
  \end{equation*}
  By the hypothesis $\kassoc{\rmS}{Y,X}{\rmZ}$ we have $\rmZ\in\mathcal{W}$,
  so $\mathcal{W}$ is nonempty; and since $\rmZ$ is finite, so is
  $\mathcal{W}$. Hence $\mathcal{W}$ contains an element $\rmZ'$ of minimum
  cardinality. Every proper subset $\rmZ^{*}\subsetneq\rmZ'$ satisfies
  $|\rmZ^{*}|<|\rmZ'|$ and $\rmZ^{*}\subseteq\rmZ$, so by the minimality of
  $\rmZ'$ we have $\rmZ^{*}\notin\mathcal{W}$, i.e.\
  $\neg\,\kassoc{\rmS}{Y,X}{\rmZ^{*}}$. Together with
  $\kassoc{\rmS}{Y,X}{\rmZ'}$, this is by \Cref{eq:k-dep-dep} exactly
  $\kdep{\rmS}{Y,X}{\rmZ'}$. As $\rmZ'\subseteq\rmZ$, this establishes
  \Cref{eq:11}.
\end{proof}

%%%% k dep to inter
\kdepinter*
\begin{proof}
  Assume $\kdep{\emptyset}{Y,X}{\rmZ}$. For the variable $X$, the
  inter-association clause is $\kassoc{\emptyset}{Y,X}{\rmZ}$, i.e.\
  $Y\dep{\pr}X\mid\rmZ$, which is exactly the first conjunct of
  \Cref{eq:k-dep-dep} and hence immediate.

  Now suppose, towards a contradiction, that the clause fails for some
  $Z\in\rmZ$, i.e.\ $Y \indep{\pr} Z \mid X \cup (\rmZ \setminus Z)$.
  Since $\rmZ\setminus Z$ is a proper subset of $\rmZ$ and the
  separating set is empty, the second conjunct of \Cref{eq:k-dep-dep}
  gives $Y \indep{\pr} X \mid \rmZ \setminus Z$. Applying the
  contraction and then weak union graphoid axioms of
  \Cref{def:graphoid},
  \begin{equation}
    \begin{aligned}
      &(Y \indep{\pr} Z \mid X \cup (\rmZ \setminus Z))
        \wedge (Y \indep{\pr} X \mid \rmZ \setminus Z)\\
      &\implies Y \indep{\pr} \{X, Z\} \mid \rmZ \setminus Z
      \implies Y \indep{\pr} X \mid \rmZ,
    \end{aligned}\label{eq:k-dep-inter-eq4}
  \end{equation}
  contradicting $Y\dep{\pr}X\mid\rmZ$. Hence
  $Y\dep{\pr}Z\mid X\cup(\rmZ\setminus Z)$ for every $Z\in\rmZ$, which
  together with $Y\dep{\pr}X\mid\rmZ$ is precisely
  $\kinter{\emptyset}{Y,\rmZ\cup X}$.
\end{proof}

%%% Proofs of Theorems
\section{Correctness of {\method}}\label{apx:proofs-thm}
\citest*

\begin{proof}
  The function \textit{find\_cond} first iterates through all the
  subsets $\rmX'\in\powerleq{k-|\rmX|}{\rmS}$ and for each $\rmX'$ where
  $Y \dep{\pre} X \mid \rmX\cup\rmX'$ (line 9 in \textit{find\_cond}), it passes
  $\rmX'$ to the function \textit{no\_seps} (line 10 in
  \textit{find\_cond}) to determine if there are any subsets
  $\rmS'\in\powerleq{l}{\rmS\setminus\rmX'}$ that results in
  $Y \indep{\pre} X \mid \rmX\cup\rmX'\cup \rmS'$ (line 4 in \textit{no\_seps}).
  As soon as \textit{no\_seps} finds such a $\rmS'$, it returns
  \texttt{False} (line 4 in \textit{no\_seps}), and \textit{find\_cond}
  will need to go to the next iteration. Otherwise, if \textit{no\_seps}
  is unable to find any such subset $\rmS'$, it returns \texttt{True}
  (line 5 in \textit{no\_seps}), and \textit{find\_cond} then
  immediately returns $\rmZ=\rmX'\cup\rmX$ (line 11 in
  \textit{find\_cond}) as there are no $l$-bounded subsets
  $\rmS'\in\powerleq{l}{\rmS\setminus\rmZ}$ that render $Y$ and $X$
  conditionally independent given
  $\rmZ\cup\rmS'$. Here we use that the external set $\rmX$ is disjoint from
  $\rmS$, so $\rmS\setminus\rmZ=\rmS\setminus\rmX'$; and since
  $\emptyset\in\powerleq{l}{\rmS\setminus\rmX'}$, the dependence check
  $Y\dep{\pre}X\mid\rmZ$ on line 9 is precisely the $\rmS'=\emptyset$
  instance tested by \textit{no\_seps}, so a returned $\rmZ$ satisfies
  \Cref{eq:25} in full. However, if \textit{find\_cond} completes all its iteration
  without returning early, then we know that all subset $\rmX'$ has some
  $l$-bounded subset $\rmS'$ that will cause
  $Y\indep{\pre}X\mid\rmX'\cup\rmX\cup\rmS'$. Therefore
  \textit{find\_cond} returns the failure value $\bot$ in this case
  (line 12 in \textit{find\_cond}).
\end{proof}

\assocmine*
\begin{proof}
  First consider a call that returns a nonempty set.
  \textit{find\_inter\_dep} does so only by returning
  \textit{make\_strict}$(Y,X,\rmZ,\rmS,k,l)$ after a call
  \textit{find\_cond}$(Y,X,\rmX,\rmS,k,l)$ has returned some
  $\rmZ\neq\bot$, for some $X\in\verts\setminus(Y\cup\rmS\cup\rmX)$
  (as $\rmS$ is passed unchanged through every recursive call, any such
  $X$ satisfies $X\in\verts\setminus(Y\cup\rmS)=\verts'$, matching the
  theorem's candidate set). By
  \Cref{thm:citest}, $\rmZ=\rmX\cup\rmX'$ for some
  $\rmX'\in\powerleq{k-|\rmX|}{\rmS}$ such that
  $Y\dep{\pre}X\mid\rmZ\cup\rmS'$ holds for all
  $\rmS'\in\powerleq{l}{\rmS\setminus\rmX'}$. As the dependence set $\rmX$ is
  disjoint from $\rmS$, we have $\rmS\setminus\rmX'=\rmS\setminus\rmZ$, so
  $Y\dep{\pre}X\mid\rmZ\cup\rmS'$ holds for all
  $\rmS'\in\powerleq{l}{\rmS\setminus\rmZ}$ with
  $\rmZ\in\powerleq{k}{\verts\setminus\{Y,X\}}$---the estimated $l$-bounded
  association claimed in the theorem. Since
  \textit{make\_strict} only removes variables from $\rmZ$, the returned
  set is $X\cup\rmZ_{R}$ for some $\rmZ_{R}\subseteq\rmZ$, as claimed.

  Next consider a call that returns the empty set. Here every internal
  \textit{find\_cond} call must have failed, so \textit{find\_inter\_dep}
  traverses the full enumeration tree of the dependence sets
  $\rmX\in\powerleq{k}{\verts\setminus(Y\cup\rmS)}$ (lines 13--17 of
  \textit{find\_inter\_dep}), and for each visits every candidate
  $X\in\verts\setminus(Y\cup\rmS\cup\rmX)$ through
  \textit{find\_cond}$(Y,X,\rmX,\rmS,k,l)$, which by \Cref{thm:citest}
  searches all $\rmX'\in\powerleq{k-|\rmX|}{\rmS}$. Every conditioning
  set $\rmZ\in\powerleq{k}{\verts\setminus\{Y,X\}}$ decomposes uniquely
  into its part outside and its part inside the candidate blanket,
  $\rmZ=(\rmZ\setminus\rmS)\cup(\rmZ\cap\rmS)=\rmX\cup\rmX'$, so this
  enumeration covers every pair $(X,\rmZ)$ with $X\in\verts'$ and
  $\rmZ\in\powerleq{k}{\verts\setminus\{Y,X\}}$. By \Cref{thm:citest},
  the call for $(X,\rmX)$ returns $\bot$ exactly when every such $\rmX'$
  admits a separator $\rmS'\in\powerleq{l}{\rmS\setminus\rmX'}$ with
  $Y\indep{\pre}X\mid\rmX\cup\rmX'\cup\rmS'$. Therefore
  \textit{find\_inter\_dep} returns the empty set if and only if
  \Cref{eq:fid-empty} holds. The distinguished value $\bot$ lets the
  callers tell such a failure from a success that returns the empty
  conditioning set.
\end{proof}

\algkomb*
\begin{proof}
  Throughout we assume that the \gls{ci} decisions are accurate, so
  $\pre$ and $\pr$ agree on every tested statement.

  In the grow phase, every nonempty set returned by
  \textit{find\_inter\_dep} contains a variable $X\notin\rmS$ by
  \Cref{thm:assocmine}, so each iteration of the repeat loop strictly
  grows $\rmS$ and the phase terminates. Upon termination
  \textit{find\_inter\_dep} returned the empty set, so by
  \Cref{thm:assocmine} the condition in \Cref{eq:fid-empty} holds in
  $\pre$, and by accuracy in $\pr$. Fix any
  $X\in\verts\setminus(\rmS\cup Y)$. For every
  $\rmZ\in\powerleq{k}{\verts\setminus\{Y,X\}}$, \Cref{eq:fid-empty}
  supplies an $\rmS'\in\powerleq{l}{\rmS\setminus\rmZ}
  \subseteq\powerleq{l}{\rmS}$ with $Y\indep{\pr}X\mid\rmZ\cup\rmS'$;
  the $l$-bounded separator assumption (\Cref{def:bound-sep}) then
  lifts this, for every such $\rmZ$, to the unbounded statement that
  some $\rmS'\subseteq\rmS$ gives $Y\indep{\pr}X\mid\rmZ\cup\rmS'$,
  which is the antecedent of $k$-order faithfulness
  (\Cref{eq:k-faith-indep}). Hence $Y\indep{\gr}X\mid\rmS$ for every
  $X\in\verts\setminus(\rmS\cup Y)$. Since the parents and children of
  $Y$---and, for an undirected $\gr$, its neighbours---can never be
  separated from $Y$, they must all lie in $\rmS$. Any spouse $X$
  shares a child $C\in\text{ch}(Y)\subseteq\rmS$, so the collider path
  $Y\rightarrow C\leftarrow X$ is active given $\rmS$; hence
  $Y\dep{\gr}X\mid\rmS$, which forces $X\in\rmS$ as well. Therefore
  $\mbs{\gr}(Y)\subseteq\rmS$, establishing \Cref{eq:27}.

  During the shrink phase we maintain the invariant
  $\mbs{\gr}(Y)\subseteq\rmS$, which holds on entry by \Cref{eq:27}.
  First, we show that no true member is ever removed. For
  $X\in\mbs{\gr}(Y)$, the in-blanket witness hypothesis of
  \Cref{thm:algkomb} supplies a set
  $\rmZ\in\powerleq{k}{\mbs{\gr}(Y)\setminus X}
  \subseteq\powerleq{k}{\rmS\setminus X}$ such that
  $Y\dep{\pr}X\mid\rmZ\cup\rmS'$ holds for all
  $\rmS'\in\powerleq{l}{(\rmS\setminus X)\setminus\rmZ}$. By accuracy and \Cref{thm:citest}
  (with candidate blanket $\rmS\setminus X$, external dependence set
  $\rmX=\emptyset$, and $\rmX'=\rmZ$), the call
  \textit{find\_cond}$(Y,X,\emptyset,\rmS\setminus X,k,l)$ does not
  return $\bot$, so $X$ is never selected for removal (a member whose
  only witness is $\rmZ=\emptyset$ makes \textit{find\_cond} return the
  empty set rather than $\bot$).

  Next, we show that every non-member is removed once selected. Let
  $X\in\rmS\setminus\mbs{\gr}(Y)$; by the invariant
  $\mbs{\gr}(Y)\subseteq\rmS\setminus X$. We first observe that
  $\mbs{\gr}(Y)$ separates $Y$ from $X$ even after conditioning on any
  further set $\rmZ\subseteq\verts\setminus\{Y,X\}$. Since
  $X\notin\mbs{\gr}(Y)$, $X$ is not adjacent to $Y$, so every path
  between them has at least one interior vertex; write such a path as
  $Y=N_{0}\sim N_{1}\sim\cdots\sim N_{m}=X$ with $m\geq 2$. Its first
  interior vertex $N_{1}$ is adjacent to $Y$ and hence lies in
  $\mbs{\gr}(Y)$. If $N_{1}$ is a non-collider on the path, it blocks
  the path. If instead $N_{1}$ is a collider
  $Y\rightarrow N_{1}\leftarrow N_{2}$, then $N_{1}$ is a common child of
  $Y$ and its successor $N_{2}$ on the path, so $N_{2}$ is a spouse of
  $Y$ and hence $N_{2}\in\mbs{\gr}(Y)$. In particular $N_{2}\neq X$ (as
  $X\notin\mbs{\gr}(Y)$), so $N_{2}$ is a genuine interior vertex; and
  since the edge $N_{2}\rightarrow N_{1}$ points out of $N_{2}$, it is a
  non-collider on the path and blocks it. For an undirected $\gr$, $N_{1}$ is instead a neighbour of
  $Y$ and blocks the path directly. As each path is thus blocked at a
  conditioned vertex of $\mbs{\gr}(Y)$ irrespective of
  $\rmZ$~\citep{pearl1988,lauritzen1996},
  $Y\indep{\gr}X\mid\mbs{\gr}(Y)\cup\rmZ$ for every
  $\rmZ\in\powerleq{k}{\verts\setminus\{Y,X\}}$.

  The Global Markov
  Property transfers this to $\pr$ with the separator
  $\rmS'\coloneq\mbs{\gr}(Y)\setminus\rmZ\subseteq\rmS\setminus X$, so
  the unbounded side of \Cref{def:bound-sep}---at the separating set
  $\rmS\setminus X$---holds for every such $\rmZ$. The biconditional
  then yields some $\rmS'\in\powerleq{l}{\rmS\setminus X}$; discarding
  any overlap with $\rmZ$ leaves $\rmZ\cup\rmS'$ unchanged and gives
  $\rmS'\in\powerleq{l}{(\rmS\setminus X)\setminus\rmZ}$ with
  $Y\indep{\pr}X\mid\rmZ\cup\rmS'$. By accuracy and \Cref{thm:citest},
  \textit{find\_cond}$(Y,X,\emptyset,\rmS\setminus X,k,l)$ returns
  $\bot$, so $X$ is removed once selected.

  Each iteration of the while loop removes one variable, so the phase
  terminates; it exits only when no $X\in\rmS$ yields $\bot$, whence
  $\rmS\setminus\mbs{\gr}(Y)=\emptyset$ and, with the invariant,
  $\rmS=\mbs{\gr}(Y)$, establishing \Cref{eq:28}.
\end{proof}

\section{Complexity of {\method}}\label{apx:proof-complex}

\kombcomplexity*
\begin{proof}
  We charge $O(N)$ to each \gls{ci} test and bound the running time by
  analysing the three sub-algorithms of \Cref{sec:algo} from the
  innermost outwards.

  First consider \textit{find\_cond} (\Cref{alg:ci-test}) for a fixed
  candidate variable $X$ and external dependence set $\rmX$ with
  $|\rmX|=i$. It iterates over every subset
  $\rmX'\in\powerleq{k-i}{\rmS}$, of which there are at most
  $\sum_{i'=0}^{k-i}\binom{K}{i'}$ since $|\rmS|\leq K$. For each $\rmX'$
  it performs one \gls{ci} test and, whenever that test indicates a
  dependence, calls \textit{no\_seps}, which iterates over every
  separator $\rmS'\in\powerleq{l}{\rmS\setminus\rmX'}$ and performs one
  further \gls{ci} test per separator. Writing $|\rmX'|=i'$, there are at
  most $\sum_{j=0}^{l}\binom{K-i'}{j}$ such separators, so a single
  call to \textit{find\_cond} costs
  \begin{equation}
    \label{eq:cost-findcond}
    O\!\left(
      \sum_{i'=0}^{k-i}\binom{K}{i'}
      \sum_{j=0}^{l}\binom{K-i'}{j}\, N
    \right),
  \end{equation}
  where the convention $\binom{n}{m}=0$ for $n<m$ accounts for the
  separator set $\rmS\setminus\rmX'$ being exhausted.

  Next consider \textit{find\_inter\_dep} (\Cref{alg:assoc-mine}). It
  explores the enumeration tree of $\verts\setminus Y$ up to subsets of
  cardinality $k$ (lines 13--17 of \textit{find\_inter\_dep}), and hence
  visits at most $\sum_{i=0}^{k}\binom{n}{i}$ dependence sets
  $\rmX\in\powerleq{k}{\verts\setminus Y}$ (a loose bound on the exact
  count $\sum_{i=0}^{k}\binom{n-1}{i}$, as $|\verts\setminus Y|=n-1$). For every such $\rmX$ it
  loops over the $O(n)$ candidate variables
  $X\in\verts\setminus(Y\cup\rmS\cup\rmX)$ and invokes
  \textit{find\_cond}. Multiplying these two factors by
  \Cref{eq:cost-findcond} bounds the cost of a single association-mining
  sweep by
  \begin{equation}
    \label{eq:cost-sweep}
    O\!\left(
      n \sum_{i=0}^{k}\binom{n}{i}
      \left[
        \sum_{i'=0}^{k-i}\binom{K}{i'}
        \sum_{j=0}^{l}\binom{K-i'}{j}\, N
      \right]
    \right).
  \end{equation}
  Each successful sweep additionally invokes \textit{make\_strict}
  exactly once; it iterates over the found dependence set $\rmZ$---of
  size at most $\min(k,n-2)$, since $\rmZ\subseteq\verts\setminus\{Y,X\}$
  and $|\rmZ|\leq k$---and performs a single \textit{no\_seps} call per
  element, at cost
  $O\!\left(\min(k,n-2)\sum_{j=0}^{l}\binom{K}{j}\,N\right)$. This is dominated by
  the per-sweep \textit{find\_cond} total already counted in
  \Cref{eq:cost-sweep}---whose bracketed factor already contains the
  term $\sum_{j=0}^{l}\binom{K}{j}$ (at $i'=0$) multiplied by the outer
  factor $n\sum_{i=0}^{k}\binom{n}{i}\geq n\geq\min(k,n-2)$---and is
  therefore absorbed into the bound.

  Finally, consider {\method} itself (\Cref{alg:komb}). Each successful
  call to \textit{find\_inter\_dep} in the grow phase returns a
  nonempty set $X\cup\rmZ$ that is added to $\rmS$, so the candidate
  Markov blanket grows by at least one variable per successful call; as
  $|\rmS|\leq K$ throughout, the grow phase performs at most $O(K)$
  association-mining sweeps. The shrink phase re-evaluates its guard after
  each removal, calling \textit{find\_cond} at most $O(K^{2})$ times over the
  at most $K$ removals; as $K\leq n-1$, this cost is dominated by that of the
  grow phase. Multiplying \Cref{eq:cost-sweep}
  by these $O(K)$ sweeps yields the bound in \Cref{eq:complexity},
  completing the proof.
\end{proof}


%%% Implementation
% \section{Implementation Details of kOMB}\label{apx:impl}
% The main conditional independence test we use in our implementation is
% the G-test. Conducting this test first involves computing the
% G-statistic,
% \begin{equation}
%   \label{eq:g-stat}
%   G = 2 \sum_{i} O_{i}\cdot \ln\Biggl(\frac{O_{i}}{E_{i}}\Biggr),
% \end{equation}
% where $O_{i}\geq 0$ are the observed counts in the cells and $E_{i} > 0$
% are the expected counts under the distribution assumed in the null
% hypothesis.

% In the implementation of {\method}, we are specifically interested in
% using the G-test to test the null hypothesis,
% \begin{equation*}
%   H_{0}(\rmS): X \indep{\pre} Y \mid \rmS
% \end{equation*}
% The G-statistic for testing this hypothesis is
% then~\citep{tsamardinos2006}
% \begin{align*}
%   G(\rmS)
%   &= 2 \sum_{x,y,\rvs}
%     O(x,y,\rvs)\cdot
%     \ln\Biggl(
%     \frac{P(x,y\mid\rvs) P(\rvs)}
%     {P(x\mid\rvs)P(y\mid\rvs)
%     P(\rvs)}\Biggr)\\
%   &= 2 \sum_{x,y,\rvs}
%     O(x,y,\rvs)\cdot
%     \ln\Biggl(
%     \frac{O(x,y,\rvs) O(\rvs)}
%     {O(x,\rvs)O(y,\rvs)}\Biggr).
% \end{align*}
% Under the null hypothesis, the distribution of this statistic $G(\rmS)$
% is then approximately $\chi^{2}(\df(\rmS))$ distributed
% where~\citep{tsamardinos2006}
% \begin{align*}
%   \df(\rmS) &= (|Y|-1) \cdot (|X| - 1)\cdot
%               \Biggl(\prod_{S\in\rmS}|S|\Biggr),
% \end{align*}
% is the degrees of freedomn in the statistical test.

% However, the G-test starts to break down when $\rmS$ grows large
% enough. The common cut-off point to say that the G-test is too weak for
% a CI test is when $\df(\rmS)\times 5$ is greater than the size of the
% available samples.

% When this occurs, we switch to using the SCI test proposed by
% \citet{marx2019}. We defer to the original paper for the details of the
% SCI test.



%%% Synthetic Bayesian Networks
\section{Boolean Functions for the Synthetic Bayesian Networks}\label{apx:bool-func}
\begin{description}[align=right,labelwidth=7em]
\item[parity:]     
  $f(X_{1},\ldots,X_{3}) = (\sum_{i=1}^{3}X_{i}) \pmod{2}$
\item[ex-$s$:] 
  $f(X_{1},\ldots,X_{3}; s) = \bigl[(\sum_{i=1}^{3}X_{i}) = s\bigr]$
\item[and:]       
  $f(X_{1},\ldots,X_{3}) = \bigl[(\sum_{i=1}^{3}X_{i}) = 3\bigr]$
\item[or:]         
  $f(X_{1},\ldots,X_{3}) = \bigl[(\sum_{i=1}^{3}X_{i}) \geq 1\bigr]$
\end{description}
where $s\in\{1,2\}$.

%%% Separator sizes: empirical support for the l-bounded separator assumption
\section{Separator Sizes in the Benchmark Datasets}\label{apx:dsep}
To gauge how restrictive the $l$-bounded separator assumption
(\Cref{def:bound-sep}) is in practice, we examine the benchmark networks
of \Cref{sec:exp-real}. We use every variable in turn as a target and,
for every pair formed by a target and a variable outside its Markov
blanket, compute a minimal d-separator drawn from the true Markov
blanket. \Cref{tab:dsep} reports the mean and worst-case (maximum)
separator size and the fraction of pairs whose separator has size at
most three. Separators are small across all four networks---the mean
never exceeds $2.19$ and size-$\leq\!3$ separators cover at least $83\%$
of pairs---and no pair had to be skipped, i.e.\ a separator within the
true Markov blanket always existed. The bound is loosest on Barley,
where the largest separator reaches size $9$; this is also the network
on which {\method} is the most expensive and least accurate
(\Cref{sec:exp-real}), illustrating that the difficulty of \gls{mb}
discovery tracks the separator sizes the assumption must accommodate.
Overall, a modest separator bound $l$ already suffices to capture most
of the conditional independencies needed to recover the Markov blanket
on these networks.

\begin{table}[h]
  \centering
  \caption{Minimal d-separator sizes restricted to the true Markov
    blanket, computed using every variable of each benchmark network as
    a target. ``Targets'' is the number of targets (all variables) and
    ``Pairs'' the number of target--non-blanket pairs; ``Mean sep.\ size''
    and ``Max'' are the mean and maximum separator size; and
    ``Cover $\leq 3$'' is the fraction whose minimal separator has size
    at most three.}
  \label{tab:dsep}
  \resizebox{\columnwidth}{!}{\begin{tabular}{llllll}
\toprule
Network & Targets & Pairs & Mean sep.\ size & Max & Cover $\leq 3$ \\
\midrule
Alarm1 & 37 & 1202 & 0.621 & 3 & 1.000 \\
Barley & 48 & 2004 & 1.739 & 9 & 0.835 \\
Insurance & 27 & 562 & 2.190 & 7 & 0.925 \\
Mildew & 35 & 1030 & 1.065 & 6 & 0.976 \\
\bottomrule
\end{tabular}
}
\end{table}

%%% Full benchmark results / Ablation study
\section{Ablation}\label{apx:ablation}
For completeness, we report the full results of the experiments in
\Cref{sec:exp-real} for every method and every $(k,l)$ configuration of
{\method}. \Cref{tab:real-full} gives the performance of all methods,
while \Cref{tab:runtime-full} gives their runtimes in seconds. In both
tables, ``--'' denotes a configuration that failed to complete within the
allotted time budget.

\Cref{tab:kl-ablation} summarises the F1 score of {\method} across all
synthetic and benchmark datasets for every $(k,l)$ configuration. F1
generally improves from $k=0$ to $k=2$---higher orders expose
dependencies that lower orders miss---but the gains flatten or reverse
beyond $k=2$ as the larger conditioning sets incur both higher runtime
(\Cref{tab:runtime-full}) and more error-prone \gls{ci} tests. Setting
$l=k$ with $k\leq 2$ therefore offers a robust default, the heuristic we
adopt in \Cref{sec:complexity}. For a fixed $k$, raising $l$ beyond $k$
does not consistently improve F1 (\Cref{tab:kl-ablation}), which
motivates simply setting $l=k$ rather than tuning $l$ separately.

\begin{table*}[h]
  \centering
  \caption{F1 score across datasets for each $(k,l)$ configuration of {\method}.}
  \label{tab:kl-ablation}
  \setlength{\tabcolsep}{5pt}
\begin{tabular}{lccccccccc}
\toprule
Method & parity & ex-1 & ex-2 & and & or & Alarm1 & Barley & Insurance & Mildew \\
\midrule

kOMB
& 0.0250 & 0.2825 & 0.2675 & 0.4100 & 0.3500 & -- & -- & -- & -- \\
{\scriptsize $(k=0,l=1)$} & {\scriptsize $\pm$0.1104} & {\scriptsize $\pm$0.3720} & {\scriptsize $\pm$0.3737} & {\scriptsize $\pm$0.4050} & {\scriptsize $\pm$0.3508} & -- & -- & -- & -- \\

kOMB
& 0.0250 & 0.2825 & 0.2675 & 0.4100 & 0.3500 & -- & -- & -- & -- \\
{\scriptsize $(k=0,l=2)$} & {\scriptsize $\pm$0.1104} & {\scriptsize $\pm$0.3720} & {\scriptsize $\pm$0.3737} & {\scriptsize $\pm$0.4050} & {\scriptsize $\pm$0.3508} & -- & -- & -- & -- \\

kOMB
& 0.0250 & 0.2825 & 0.2675 & 0.4100 & 0.3500 & -- & -- & -- & -- \\
{\scriptsize $(k=0,l=3)$} & {\scriptsize $\pm$0.1104} & {\scriptsize $\pm$0.3720} & {\scriptsize $\pm$0.3737} & {\scriptsize $\pm$0.4050} & {\scriptsize $\pm$0.3508} & -- & -- & -- & -- \\

kOMB
& 0.0500 & 0.6750 & 0.6750 & \textbf{0.6175} & 0.5600 & \textbf{0.7283} & \textbf{0.4323} & 0.7155 & \textbf{0.5035} \\
{\scriptsize $(k=1,l=1)$} & {\scriptsize $\pm$0.2207} & {\scriptsize $\pm$0.4743} & {\scriptsize $\pm$0.4743} & {\scriptsize $\pm$0.3741} & {\scriptsize $\pm$0.4349} & {\scriptsize $\pm$0.1630} & {\scriptsize $\pm$0.1385} & {\scriptsize $\pm$0.1322} & {\scriptsize $\pm$0.1637} \\

kOMB
& 0.0500 & 0.6750 & 0.6750 & \textbf{0.6175} & 0.5600 & 0.7245 & 0.3510 & \textbf{0.7377} & 0.4686 \\
{\scriptsize $(k=1,l=2)$} & {\scriptsize $\pm$0.2207} & {\scriptsize $\pm$0.4743} & {\scriptsize $\pm$0.4743} & {\scriptsize $\pm$0.3741} & {\scriptsize $\pm$0.4349} & {\scriptsize $\pm$0.1983} & {\scriptsize $\pm$0.1344} & {\scriptsize $\pm$0.1594} & {\scriptsize $\pm$0.1255} \\

kOMB
& 0.0500 & 0.6750 & 0.6750 & \textbf{0.6175} & 0.5600 & 0.6915 & 0.3510 & 0.6728 & 0.4686 \\
{\scriptsize $(k=1,l=3)$} & {\scriptsize $\pm$0.2207} & {\scriptsize $\pm$0.4743} & {\scriptsize $\pm$0.4743} & {\scriptsize $\pm$0.3741} & {\scriptsize $\pm$0.4349} & {\scriptsize $\pm$0.1825} & {\scriptsize $\pm$0.1344} & {\scriptsize $\pm$0.1092} & {\scriptsize $\pm$0.1255} \\

kOMB
& \textbf{1.0000} & \textbf{1.0000} & \textbf{1.0000} & \textbf{0.6175} & \textbf{0.7100} & 0.7280 & 0.3590 & 0.7309 & 0.4737 \\
{\scriptsize $(k=2,l=1)$} & {\scriptsize $\pm$0.0000} & {\scriptsize $\pm$0.0000} & {\scriptsize $\pm$0.0000} & {\scriptsize $\pm$0.3741} & {\scriptsize $\pm$0.3842} & {\scriptsize $\pm$0.1788} & {\scriptsize $\pm$0.1481} & {\scriptsize $\pm$0.1388} & {\scriptsize $\pm$0.1318} \\

kOMB
& \textbf{1.0000} & \textbf{1.0000} & \textbf{1.0000} & \textbf{0.6175} & \textbf{0.7100} & 0.6903 & 0.3510 & 0.7127 & 0.4686 \\
{\scriptsize $(k=2,l=2)$} & {\scriptsize $\pm$0.0000} & {\scriptsize $\pm$0.0000} & {\scriptsize $\pm$0.0000} & {\scriptsize $\pm$0.3741} & {\scriptsize $\pm$0.3842} & {\scriptsize $\pm$0.1831} & {\scriptsize $\pm$0.1344} & {\scriptsize $\pm$0.1473} & {\scriptsize $\pm$0.1255} \\

kOMB
& \textbf{1.0000} & \textbf{1.0000} & \textbf{1.0000} & \textbf{0.6175} & \textbf{0.7100} & 0.6915 & 0.3510 & 0.6538 & 0.4686 \\
{\scriptsize $(k=2,l=3)$} & {\scriptsize $\pm$0.0000} & {\scriptsize $\pm$0.0000} & {\scriptsize $\pm$0.0000} & {\scriptsize $\pm$0.3741} & {\scriptsize $\pm$0.3842} & {\scriptsize $\pm$0.1825} & {\scriptsize $\pm$0.1344} & {\scriptsize $\pm$0.0862} & {\scriptsize $\pm$0.1255} \\

kOMB
& -- & -- & -- & -- & -- & 0.6834 & -- & 0.6991 & 0.4686 \\
{\scriptsize $(k=3,l=1)$} & -- & -- & -- & -- & -- & {\scriptsize $\pm$0.1741} & -- & {\scriptsize $\pm$0.1375} & {\scriptsize $\pm$0.1255} \\

kOMB
& -- & -- & -- & -- & -- & 0.6915 & -- & 0.6538 & 0.4686 \\
{\scriptsize $(k=3,l=2)$} & -- & -- & -- & -- & -- & {\scriptsize $\pm$0.1825} & -- & {\scriptsize $\pm$0.0862} & {\scriptsize $\pm$0.1255} \\

kOMB
& -- & -- & -- & -- & -- & 0.6915 & -- & 0.6516 & 0.4686 \\
{\scriptsize $(k=3,l=3)$} & -- & -- & -- & -- & -- & {\scriptsize $\pm$0.1825} & -- & {\scriptsize $\pm$0.0828} & {\scriptsize $\pm$0.1255} \\
\bottomrule
\end{tabular}

\end{table*}

\begin{table}[h]
  \centering
  \caption{Full results on benchmark datasets.}
  \label{tab:real-full}
  \setlength{\tabcolsep}{5pt}
\begin{tabular}{lcccc}
\toprule
Method & Alarm1 & Barley & Insurance & Mildew \\
\midrule

BAMB
& 0.6742 & 0.3385 & 0.6588 & 0.4945 \\
& {\scriptsize $\pm$0.1771} & {\scriptsize $\pm$0.1226} & {\scriptsize $\pm$0.0797} & {\scriptsize $\pm$0.1084} \\

GS
& 0.3434 & 0.2113 & 0.4615 & 0.2867 \\
& {\scriptsize $\pm$0.1503} & {\scriptsize $\pm$0.0807} & {\scriptsize $\pm$0.1477} & {\scriptsize $\pm$0.1581} \\

HITON\_MB
& 0.7636 & 0.3400 & 0.6746 & 0.4918 \\
& {\scriptsize $\pm$0.1861} & {\scriptsize $\pm$0.1284} & {\scriptsize $\pm$0.1594} & {\scriptsize $\pm$0.1487} \\

IAMB
& 0.6931 & 0.3373 & 0.6293 & 0.4677 \\
& {\scriptsize $\pm$0.1827} & {\scriptsize $\pm$0.1223} & {\scriptsize $\pm$0.0778} & {\scriptsize $\pm$0.1224} \\

LRH
& 0.6813 & 0.3373 & 0.6281 & 0.4674 \\
& {\scriptsize $\pm$0.1777} & {\scriptsize $\pm$0.1223} & {\scriptsize $\pm$0.0782} & {\scriptsize $\pm$0.1230} \\

MMMB
& 0.7692 & 0.3298 & 0.6978 & 0.4861 \\
& {\scriptsize $\pm$0.1874} & {\scriptsize $\pm$0.1274} & {\scriptsize $\pm$0.1405} & {\scriptsize $\pm$0.1810} \\

PCMB
& 0.7371 & 0.2284 & 0.6299 & 0.4156 \\
& {\scriptsize $\pm$0.2127} & {\scriptsize $\pm$0.1199} & {\scriptsize $\pm$0.1291} & {\scriptsize $\pm$0.1445} \\

STMB
& 0.6598 & 0.2145 & 0.5648 & 0.2734 \\
& {\scriptsize $\pm$0.1627} & {\scriptsize $\pm$0.1281} & {\scriptsize $\pm$0.1324} & {\scriptsize $\pm$0.1765} \\

kOMB
& 0.7804 & \textbf{0.5158} & 0.7136 & \textbf{0.6220} \\
{\scriptsize $(k=1,l=1)$} & {\scriptsize $\pm$0.1145} & {\scriptsize $\pm$0.1964} & {\scriptsize $\pm$0.1328} & {\scriptsize $\pm$0.1595} \\

kOMB
& 0.8187 & 0.3128 & \textbf{0.7336} & 0.4360 \\
{\scriptsize $(k=1,l=2)$} & {\scriptsize $\pm$0.1944} & {\scriptsize $\pm$0.1582} & {\scriptsize $\pm$0.1624} & {\scriptsize $\pm$0.1394} \\

kOMB
& 0.7923 & 0.2498 & 0.7076 & 0.3726 \\
{\scriptsize $(k=1,l=3)$} & {\scriptsize $\pm$0.1885} & {\scriptsize $\pm$0.1140} & {\scriptsize $\pm$0.1304} & {\scriptsize $\pm$0.1558} \\

kOMB
& 0.7775 & 0.4588 & 0.7249 & 0.5059 \\
{\scriptsize $(k=2,l=1)$} & {\scriptsize $\pm$0.1511} & {\scriptsize $\pm$0.1825} & {\scriptsize $\pm$0.1464} & {\scriptsize $\pm$0.1458} \\

kOMB
& \textbf{0.8209} & 0.3128 & \textbf{0.7336} & 0.4360 \\
{\scriptsize $(k=2,l=2)$} & {\scriptsize $\pm$0.1914} & {\scriptsize $\pm$0.1582} & {\scriptsize $\pm$0.1624} & {\scriptsize $\pm$0.1394} \\

kOMB
& 0.7967 & 0.2498 & 0.7058 & 0.3726 \\
{\scriptsize $(k=2,l=3)$} & {\scriptsize $\pm$0.1828} & {\scriptsize $\pm$0.1140} & {\scriptsize $\pm$0.1311} & {\scriptsize $\pm$0.1558} \\

kOMB
& 0.7721 & 0.0000 & 0.7227 & 0.0419 \\
{\scriptsize $(k=3,l=1)$} & {\scriptsize $\pm$0.1599} & {\scriptsize $\pm$0.0000} & {\scriptsize $\pm$0.1552} & {\scriptsize $\pm$0.1458} \\

kOMB
& 0.8001 & 0.0000 & 0.7288 & 0.3607 \\
{\scriptsize $(k=3,l=2)$} & {\scriptsize $\pm$0.1948} & {\scriptsize $\pm$0.0000} & {\scriptsize $\pm$0.1471} & {\scriptsize $\pm$0.1964} \\

kOMB
& 0.7967 & 0.0000 & 0.7058 & 0.3726 \\
{\scriptsize $(k=3,l=3)$} & {\scriptsize $\pm$0.1828} & {\scriptsize $\pm$0.0000} & {\scriptsize $\pm$0.1311} & {\scriptsize $\pm$0.1558} \\
\bottomrule
\end{tabular}

\end{table}

\begin{table}[h]
  \centering
  \caption{Full runtimes (in seconds) on benchmark datasets.}
  \label{tab:runtime-full}
  \setlength{\tabcolsep}{5pt}
\begin{tabular}{lcccc}
\toprule
Method & Alarm1 & Barley & Insurance & Mildew \\
\midrule

BAMB
& 7.439 & 2.232 & 10.930 & 1.604 \\
& {\scriptsize $\pm$9.063} & {\scriptsize $\pm$1.347} & {\scriptsize $\pm$10.543} & {\scriptsize $\pm$0.821} \\

GSMB
& 10.077 & 3.078 & 9.066 & 2.136 \\
& {\scriptsize $\pm$10.783} & {\scriptsize $\pm$1.195} & {\scriptsize $\pm$7.876} & {\scriptsize $\pm$0.700} \\

HITON\_MB
& 5.471 & 1.972 & 8.424 & 1.325 \\
& {\scriptsize $\pm$2.768} & {\scriptsize $\pm$0.561} & {\scriptsize $\pm$6.076} & {\scriptsize $\pm$0.483} \\

IAMB
& 5.858 & 2.164 & 3.367 & 2.000 \\
& {\scriptsize $\pm$4.029} & {\scriptsize $\pm$0.966} & {\scriptsize $\pm$2.024} & {\scriptsize $\pm$0.843} \\

LRH
& 178.162 & 8.256 & 92.170 & 3.820 \\
& {\scriptsize $\pm$183.927} & {\scriptsize $\pm$4.056} & {\scriptsize $\pm$51.496} & {\scriptsize $\pm$1.390} \\

MMMB
& 5.608 & 2.048 & 6.750 & 1.228 \\
& {\scriptsize $\pm$2.573} & {\scriptsize $\pm$0.697} & {\scriptsize $\pm$4.911} & {\scriptsize $\pm$0.294} \\

PCMB
& 30.272 & 9.135 & 52.643 & 6.305 \\
& {\scriptsize $\pm$20.025} & {\scriptsize $\pm$3.488} & {\scriptsize $\pm$42.369} & {\scriptsize $\pm$3.100} \\

STMB
& 10.281 & 3.232 & 34.573 & 1.233 \\
& {\scriptsize $\pm$11.339} & {\scriptsize $\pm$4.158} & {\scriptsize $\pm$80.556} & {\scriptsize $\pm$1.174} \\

kOMB
& 39.762 & 134.674 & 23.411 & 30.449 \\
{\scriptsize $(k=1,l=1)$} & {\scriptsize $\pm$28.094} & {\scriptsize $\pm$68.199} & {\scriptsize $\pm$13.446} & {\scriptsize $\pm$18.624} \\

kOMB
& 13.281 & 98.377 & 9.321 & 17.496 \\
{\scriptsize $(k=1,l=2)$} & {\scriptsize $\pm$6.715} & {\scriptsize $\pm$76.399} & {\scriptsize $\pm$5.293} & {\scriptsize $\pm$15.759} \\

kOMB
& 277.901 & -- & 180.066 & 308.468 \\
{\scriptsize $(k=2,l=1)$} & {\scriptsize $\pm$157.455} & -- & {\scriptsize $\pm$66.437} & {\scriptsize $\pm$429.181} \\

kOMB
& 206.537 & 104.350 & 102.564 & 46.457 \\
{\scriptsize $(k=2,l=2)$} & {\scriptsize $\pm$54.885} & {\scriptsize $\pm$94.693} & {\scriptsize $\pm$44.411} & {\scriptsize $\pm$76.392} \\

kOMB
& -- & -- & 1647.257 & -- \\
{\scriptsize $(k=3,l=1)$} & -- & -- & {\scriptsize $\pm$697.069} & -- \\

kOMB
& -- & 117.179 & -- & 49.248 \\
{\scriptsize $(k=3,l=2)$} & -- & {\scriptsize $\pm$119.477} & -- & {\scriptsize $\pm$83.115} \\

kOMB
& -- & -- & -- & 4.008 \\
{\scriptsize $(k=3,l=3)$} & -- & -- & -- & {\scriptsize $\pm$2.674} \\
\bottomrule
\end{tabular}

\end{table}
%%% Local Variables:
%%% mode: latex
%%% TeX-master: "main"
%%% End:
